\documentclass[11pt,a4paper]{article}
\usepackage[margin=23mm,headheight=14pt]{geometry}
\usepackage[T1]{fontenc}
\usepackage[utf8]{inputenc}
\usepackage{lmodern,microtype}
\usepackage{amsmath,amssymb,amsthm,mathtools}
\usepackage{booktabs,array,longtable,calc}
\usepackage{xcolor,fancyhdr,fvextra}
\usepackage[unicode,colorlinks=true,linkcolor=blue!40!black,urlcolor=blue!50!black]{hyperref}
\usepackage{enumitem}
\setlist{nosep,leftmargin=*}
\setlength{\parskip}{5pt}
\setlength{\parindent}{0pt}
\emergencystretch=4em
\setcounter{secnumdepth}{0}
\setcounter{tocdepth}{2}
\providecommand{\tightlist}{\setlength{\itemsep}{0pt}\setlength{\parskip}{0pt}}
\DefineVerbatimEnvironment{verbatim}{Verbatim}{fontsize=\small,breaklines=true,breakanywhere=true}
\pagestyle{fancy}
\fancyhf{}
\lhead{\small Erd\H{o}s 708: the seven-source case}
\rhead{\small $g(7)=9$}
\cfoot{\thepage}
\hypersetup{pdftitle={Erdos 708: g(7)=9},pdfsubject={Computer-assisted proof and exact verification record}}
\begin{document}
\thispagestyle{empty}
\vspace*{10mm}
{\Large\sffamily Erd\H{o}s--Sur\'anyi interval-product covering}\par
\vspace{6mm}
{\Huge\bfseries $\boldsymbol{g(7)=9}$}\par
\vspace{4mm}
{\large A computer-assisted proof with reproducible integer certificates}\par
\vspace{4mm}
24 November 2026\par
\vspace{7mm}
\hrule
\vspace{6mm}
\textbf{Theorem.} For any seven distinct integers $A\subseteq\mathbb Z_{\ge2}$,
put $M=\max A$. Every interval $I$ of $M$ consecutive positive integers
contains a subset $B$ of at most nine distinct elements satisfying
\[
       \prod_{a\in A}a\ \bigm|\ \prod_{b\in B}b.
\]
There are explicit $A,I$ for which eight elements never suffice.
Therefore the sharp universal bound, in the at-most convention, is nine.

\textbf{Proof architecture.} An elementary witness establishes the lower
bound. For the upper bound, a hypothetical counterexample is normalized
by its maximum and source sum. Structural lemmas leave at most three
active primes above seven. An analytic argument excludes every
$M\ge24697$. The remaining finite cases are exhausted by exact integer
calculations, with certificates valid for every interval shift.

\begin{center}
\begin{tabular}{p{32mm}p{82mm}l}
\toprule
Finite branch & Main verification record & Result\\
\midrule
No large prime & 3,413,772 source sets; 318 complete residue checks & PASS\\
One large prime & 7,006,656 source sets; 252,175 slot certificates & PASS\\
Two large primes & 116,394 final sets; 24,511 slot certificates & PASS\\
Three large primes & 850,366 core classes; all residual families certified & PASS\\
\bottomrule
\end{tabular}
\end{center}

\textbf{Dependencies and scope.} The proof uses the classical small cases
and Tripodo's stated theorems $g(4)=5$, $g(5)=6$. A separate verification of
$g(6)\le8$ is included. The new seven-source argument is a computer-assisted
proof; its certificates use exact arithmetic and the supplied verifiers
have been replayed. It is not a new Lean formalization.

\textbf{Companion archive.} The complete argument, detailed branch
appendices, source code, explicit certificate lists, independent audit
notes, and successful replay logs are included in
\nolinkurl{erdos708_g7_proof_bundle.zip}. Run:
\begin{verbatim}
python3 g7_exact/verify_all.py
\end{verbatim}
Only Python 3 and a C++17 compiler are required. Assertions must remain
enabled. The reference complete replay returned PASS in approximately
85 seconds on the verification machine.
\newpage
\tableofcontents
\newpage

Computer-assisted proof and exact verification record, 24 November 2026.

\section{Theorem and validation
status}\label{theorem-and-validation-status}

\textbf{Theorem. g(7)=9.} The explicit lower bound, normalization,
structural reductions, infinite tail, and all four finite prime-support
cases are complete and independently checked. The upper bound is
computer-assisted: analytic reductions leave a finite collection of
source configurations, and exact integer verifiers establish a universal
cover for every remaining demand. All finite branches have been
regenerated from the supplied source and their certificate checks have
returned PASS.

The argument uses the earlier small-n bounds listed in Section 2. It
does not prove the asymptotic Erdős Problem 708 or determine g(8) or
g(9).

\section{1. Definition and target}\label{definition-and-target}

For a finite set A of distinct integers at least two, put M=max(A). For
x\textgreater=0, let I=\{x+1,\ldots,x+M\}, and define

\begin{verbatim}
tau(A,I)=min{|B| : B subset I and product(A) divides product(B)}.
\end{verbatim}

The Erdős--Surányi function is

\begin{verbatim}
g(n)=sup tau(A,I),
\end{verbatim}

where the supremum ranges over all n-element source sets A and all
nonnegative integer shifts x. Cardinality is interpreted as \textbf{at
most}, not as an obligation to select exactly the same number from every
interval.

To prove g(7)=9, two separate statements are needed: one explicit
instance requiring nine points, and a universal nine-point upper bound
for every source maximum and every shift. The finite verification below
is about universal covering certificates, not a bounded sample of
shifts.

\section{2. Prior inputs and basic
divisibility}\label{prior-inputs-and-basic-divisibility}

The numerical inputs are

\begin{longtable}[]{@{}lr@{}}
\toprule\noalign{}
k & Upper bound used \\
\midrule\noalign{}
\endhead
\bottomrule\noalign{}
\endlastfoot
1 & 1 \\
2 & 2 \\
3 & 4 \\
4 & 5 \\
5 & 6 \\
6 & 8 \\
\end{longtable}

The first two are elementary. The three-source bound is the classical
Erdős--Surányi result of 1959. The external four- and five-source input
is Theorem 1.6 of John Tripodo, \emph{Sharp finite cases of the
Erdős--Surányi interval-product problem: g(4)=5 and g(5)=6}, 19 November
2026, manuscript version 1.8,
\href{https://hexagonmath.org/pdf/2611.00694v1}{primary manuscript}.
That manuscript also states the classical three-source result. It
reports a Lean registration, PALOMAR-2026-11-19-000019; the external
formalization was not rebuilt in this verification.

The six-source bound has a separate analytic and exact computational
proof in
\href{dependencies/G6_UPPER_BOUND.md}{\nolinkurl{dependencies/G6_UPPER_BOUND.md}},
using only the stated g(5)\textless=6 input. Its independent verifier is
\nolinkurl{../g7_theory/check_g6_independent.py}, with slot data in
\nolinkurl{../smooth_proof/frontier_slots6.json}. The replay checks all
1,045,917 remaining six-source subsets and returns PASS. Thus
g(6)\textless=8 is established independently of the present seven-source
argument. In particular, this proof does not rely on the faulty large-M
passage of the previously discussed Bergström preprint.

If a subset has maximum smaller than M, any suitable shorter subinterval
of I allows the corresponding g(k) bound to be used inside I.

Distinctness gives

\begin{verbatim}
product(A) divides M!.
\end{verbatim}

Also product(I)/M! is the integer binomial coefficient binomial(x+M,M).
Therefore

\begin{verbatim}
product(A) divides M! divides product(I).
\end{verbatim}

In particular, the complete interval always has sufficient valuation at
every source prime. If a selected product is missing only one unit of a
prime p, an unselected p-divisible point exists and repairs that
deficiency. More generally, missing prime units can be repaired one at a
time within their total number.

\section{3. An explicit lower bound of
nine}\label{an-explicit-lower-bound-of-nine}

Take

\begin{verbatim}
A={289,323,361,391,437,493,551}
 ={17^2,17*19,19^2,17*23,19*23,17*29,19*29}.
\end{verbatim}

Its prime demands at (17,19,23,29) are (5,5,2,2), and M=551. Set

\begin{verbatim}
N=17^2*19^2*23*29=69,587,443,
I={N-275,...,N+275}.
\end{verbatim}

The center has valuations (2,2,1,1). Every off-center point is divisible
by at most one active prime: otherwise the product of two active primes
would divide a nonzero offset of absolute value at most 275, while the
smallest such product is 323. Off-center 17- and 19-valuations are at
most one, because 17\^{}2=289 and 19\^{}2=361 also exceed 275.

If the center is chosen, the remaining demands (3,3,1,1) require eight
separate off-center points. If the center is omitted, the first two
primes require five points each and the other two require at least one
each, totaling at least twelve. Thus tau(A,I)\textgreater=9.

The nine offsets

\begin{verbatim}
0,17,34,51,19,38,57,23,29
\end{verbatim}

give a matching cover, so this instance has tau(A,I)=9 exactly.
Consequently g(7)\textgreater=9. \nolinkurl{LOWER9_SIMPLE.md} gives the
full argument, and \nolinkurl{verify_lower9_simple.py} checks every
interval point's stated restrictions and the exact product divisibility
of the displayed cover.

\section{4. Normalize a hypothetical
counterexample}\label{normalize-a-hypothetical-counterexample}

Suppose the upper bound g(7)\textless=9 fails. Choose a counterexample
with least M, and then a source set of largest sum among counterexamples
at that M and its target interval. Cases M\textless=9 are already
covered by the complete interval, so M\textgreater=10.

\subsection{4.1. The source gcd is one}\label{the-source-gcd-is-one}

If a prime p divides all seven sources, then p divides M. The target
p-multiples, after division by p, form an interval of exactly M/p
consecutive positive integers. The seven divided sources are distinct
and have maximum M/p.~If none equals one, minimality of M supplies a
nine-point cover. If one equals one, cover the other six in at most
eight points using g(6)\textless=8.

Pad the quotient cover to at least seven points when necessary, then
multiply the selected points by p.~There are at least seven quotient
positions, since seven distinct positive quotient sources exist. Padding
stays within nine points, and supplies the required additional p\^{}7.
This contradicts the counterexample. Hence gcd(A)=1.

\subsection{4.2. Upward closure and a minimum source
bound}\label{upward-closure-and-a-minimum-source-bound}

If a is a source, ka\textless=M, and ka is not a source, replace a by
ka. The new set still has seven distinct members and maximum M, its
product is a multiple of the original product, and its sum is larger. It
would remain a counterexample, contradicting maximality. Therefore A is
closed under all integer multiples at most M.

Every source consequently exceeds M/7. Otherwise its first seven
multiples exhaust the seven-element source set, giving a nontrivial
common divisor, or an eighth multiple would require too many sources.
This sharper M/7 bound uses both upward closure and gcd one.

\section{5. Elementary covering tools}\label{elementary-covering-tools}

\subsection{5.1. Cumulative-GCD slots}\label{cumulative-gcd-slots}

Let h1,\ldots,hL be positive integers at most M, with L\textless=M, such
that

\begin{verbatim}
gcd(h1*...*h_{j-1},hj) <= floor(M/j)
\end{verbatim}

for every j\textgreater=2. Then L distinct interval points have product
divisible by h1\emph{\ldots{}}hL.

Inductively suppose earlier selected points cover
H=h1\emph{\ldots{}}h\_\{j-1\}. If a fresh multiple of hj exists, append
it. Otherwise hj divides the product of the selected points, as does H,
so lcm(H,hj) divides that product. There are at least j multiples of
gcd(H,hj) in the interval and only j-1 selected points. Appending a
fresh such multiple supplies H*hj. The first step uses any multiple of
h1.

This lemma makes a finite list of small integers a universal-shift
certificate. Verification requires only integer ranges, cumulative gcds,
and final divisibility.

\subsection{5.2. Removing two sources}\label{removing-two-sources}

Cover five sources in at most six points. If the product ab of the two
removed sources factors into three integers bounded respectively by

\begin{verbatim}
floor(M/7), floor(M/8), floor(M/9),
\end{verbatim}

fresh multiples of these factors finish the cover in nine points.

A useful special case removes one occurrence of a prime p from a and one
occurrence of a prime q from b. If p\textgreater=7, q\textgreater=8, and
pq\textless=M/9, append fresh multiples of a/p, b/q, and pq. This also
works for p=q and for sources containing higher prime powers.

\subsection{5.3. A prime appearing in d
sources}\label{a-prime-appearing-in-d-sources}

Suppose a prime p divides exactly d sources, with
1\textless=d\textless7, and g(7-d)\textless=r. Write the affected
sources p\^{}\{e\_i\}c\_i, with p not dividing c\_i. If

\begin{verbatim}
p >= r+2d-1,
\end{verbatim}

then r+2d-1 points suffice.

Cover the p-free sources in r points. Select a target multiple z of one
full p-source, and at most d p-divisible points, including z, that
supply the total p-demand. Such a set exists by the nested counts of
multiples of p\^{}e: the number of sources divisible by p\^{}e is at
most floor(M/p\^{}e). A matching can include any prescribed z by
replacing the representative assigned to that source if z is not already
selected.

If z is new, its cofactor is supplied separately from the old cover, and
append the other d-1 cofactors. If z is old, the p-bank adds at most d-1
points, and append all d cofactors. Each cofactor is at most M/p, so
enough fresh multiples exist. Both cases use r+2d-1 points.

If d\textgreater=3, all affected source valuations are one,
p\^{}2\textless=M, and p\textgreater=r+2d-2, the bound improves to
r+2d-2. The p-demand p\^{}d can then be supplied in d-1 points including
z, using a square multiple and ordinary p-multiples.

\subsection{5.4. Compression of a prime
group}\label{compression-of-a-prime-group}

For the same p-sources, choose one source p\^{}\{e\_\emph{\}c\_} and put
H=product of the other d-1 cofactors. If H\textless=M and H is different
from every p-free source, cover those p-free sources together with H in
r\textless=g(8-d) points. If H=1, omit it. A p-bank of d points
including a multiple of the chosen full source, with the same overlap
accounting, completes the cover in r+d points when p\textgreater=r+d.

For d=3, r\textless=6, so an unblocked cofactor compression gives nine
points. A substantial part of the three-prime proof analyzes the cases
where all available H coincide with old sources.

\section{6. Large-prime structure}\label{large-prime-structure}

Call primes at least 11 large. All other source factors are 7-smooth.

The support lemma with d=1 and d=2 gives respectively 8+1=9 and 6+3=9.
Thus each active large prime divides at least three sources. Applying
the two-source peeling lemma to two sources divisible by p gives

\begin{verbatim}
p^2 > M/9.
\end{verbatim}

Any three large-prime occurrences, including repetitions, would have
product greater than (M/9)*11\textgreater M. Hence each source contains
at most two large-prime occurrences. There are at most four distinct
large primes, since their total source incidence is at most fourteen.

\subsection{6.1. Core-bank counting}\label{core-bank-counting}

For a chosen large-prime set, the full large part of a source is 1, p,
p\^{}2, or pq. Group sources by their nontrivial core d.~If the group
has m\_d members, choose m\_d distinct target multiples of d; these
exist because m\_d\textless=floor(M/d). A group can be required to
contain a specified multiple of one full source from that group.

At a target point shared by several groups, draw their ordinary pair
cores as edges, their squares as loops, and their singleton cores as
singleton types. If e is the number of ordinary edges, v the number of
present vertices, l the number of loops, and s the number of singleton
types, then k=e+l+s groups use the point. Their assigned valuation is
2e+2l+s, while the point supplies at least v+l. The valuation deficit is
at most 2e+l+s-v. Relative to the k assigned representatives, union size
plus unit repair has excess at most

\begin{verbatim}
1+e-v.
\end{verbatim}

With one or two primes this is never positive. With three primes it is
positive only for a triangle, and equals one there. Since
pqr\textgreater M, at most one target point can carry that triangle.
Therefore m nontrivial sources have their large demand covered in at
most m points for at most two primes, or m+1 for three primes. A
prescribed full-source point is retained.

A pair group with spare capacity can avoid the unique potential triangle
point, improving m+1 to m. If a full-source point is prescribed in
another group, it is still retained.

\subsection{6.2. Four large primes are
impossible}\label{four-large-primes-are-impossible}

Call a large prime low if it is at most sqrt(M), and high otherwise.
Each high prime has demand at least three. Each low prime has demand at
least four: demand three would mean three simple source occurrences,
which the square-gain lemma closes in nine points.

Four primes with l low ones require at least 12+l valuation units, but
at most fourteen exist, so l\textless=2. Three high primes would require
nine distinct sources, because no source can contain two high primes.
Thus exactly two are low, p1,p2, and two high, q1,q2. All inequalities
are equalities: demands are (4,4,3,3), and all seven sources have large
degree two.

Six sources have cross cores p\_i\emph{q\_j. The seventh has core p1}p2,
p1\^{}2, or p2\^{}2. Since p\_i\textgreater sqrt(M)/3 and
q\_j\textgreater sqrt(M), every cross core exceeds M/3, so its cofactor
is 1 or 2 and its multiplicity is at most two. The cross multiplicity
matrix is, up to labels,

\begin{verbatim}
[[1,2],[2,1]]   or   [[1,1],[2,2]].
\end{verbatim}

Core-bank counting covers the large demand in at most eight points. In
the loop case only the four-cycle can have positive excess, and it
occurs at at most one point. In the edge case all positive cycles
coalesce, because distinct cycles would share a single-member group.
Excess at most one gives eight points. Excess two means all five core
types meet at one point, leaving at most three selected points and
deficiencies at most (2,2,1,1). The two low-prime squares repair their
two-unit deficits in one point each, giving at most seven overall.

Require a selected point to cover one full cross source with cofactor 2;
such a source exists in a multiplicity-two group. The six cross
cofactors have product at most 16. If the seventh source is d0*s0, the
uncredited small product is at most

\begin{verbatim}
8*s0 <= 8M/d0 <= 8M/121 < M/9.
\end{verbatim}

One fresh multiple finishes the cover. Therefore a counterexample has at
most three large primes, for every M. Full details and an independent
audit are recorded in \nolinkurl{g89_theory/core_reductions.md}.

\subsection{6.3. Coprime merging of large and small
covers}\label{coprime-merging-of-large-and-small-covers}

Large-prime and small-prime demands are coprime. Independent covers can
therefore be united even if their target sets intersect. If one
prescribed point z is credited with a small cofactor c\_*, a cover of
the remaining small product need avoid only z, not every point of the
large-core cover.

This distinction is essential. It applies to a genuine large-prime core,
whose primes are all at least 11. It cannot automatically be used when a
composite 7-smooth factor has been introduced as an auxiliary atom; the
four-atom proof handles its remaining multiplier separately.

\section{7. The infinite tail:
M\textgreater=24697}\label{the-infinite-tail-m24697}

The complete proof is \nolinkurl{root/LARGE_M_CUTOFF.md}; the argument is
summarized here to show why only a finite range remains.

\subsection{7.1. Greedy packing bounds}\label{greedy-packing-bounds}

After covering five sources in six points, put any already assigned
large factors into bins M/7,M/8,M/9, and insert all remaining prime
factors, at most seven, in descending order. Failure at a prime
t\textless=7 would give

\begin{verbatim}
ab > M^3/(504*t^2) >= M^3/24696.
\end{verbatim}

Since ab\textless=M\^{}2, this is impossible when M\textgreater=24697.

A 7-smooth product Q\textless=6\^{}7 also packs into two bins M/8,M/9 in
this range. Failure would give
Q\textgreater M\textsuperscript{2/(72t)\textgreater=M}2/504, while
504*6\textsuperscript{7\textless24697}2.

\subsection{7.2. A large prime above M/9}\label{a-large-prime-above-m9}

Such a prime p has source support d in \{3,4,5,6\}; gcd one excludes
d=7. Because p\textgreater=11 and M\textless9p\textless p\^{}2, its
sources are p*c\_i with distinct c\_i in \{1,\ldots,8\}. Cover the other
sources and a p-bank, retaining a full source with largest cofactor. The
budgets are

\begin{longtable}[]{@{}lrr@{}}
\toprule\noalign{}
d & old-source bound r & r+d \\
\midrule\noalign{}
\endhead
\bottomrule\noalign{}
\endlastfoot
3 & 5 & 8 \\
4 & 4 & 8 \\
5 & 2 & 7 \\
6 & 1 & 7 \\
\end{longtable}

For d=3,4, the remaining cofactors have product at most 210 and take one
fresh point. For d=5,6, split them into two products bounded by 210 and
42 and take two fresh points. All these products are below M/9 in the
stated range. The total is at most nine.

\subsection{7.3. All large primes at most
M/9}\label{all-large-primes-at-most-m9}

If a source has large degree at most one, pair it with another source.
At most three large occurrences need placement, one in each of the three
bins; the remaining small factors pack by 7.1. Thus every source in a
counterexample has degree two.

If any core is at most M/7, put that core in the first bin and the two
individual large occurrences of another source in the other bins. Again
the small factors finish the pair. Hence every core exceeds M/7 and
every small cofactor is at most six.

If the core graph has no triangle, or a triangle edge group has spare
capacity, the core cover uses at most seven points. Two small-factor
bins cover the full small product, at most 6\^{}7.

Otherwise all three edge groups are present and saturated: their source
multiplicities equal floor(M/d), so distinctness forces their cofactors
to be exactly 1,\ldots,m\_d.~Prescribe a full source having the largest
small cofactor. The core cover uses at most eight points. If t is the
number of edge-core sources, the remaining small product is bounded by

\begin{longtable}[]{@{}lr@{}}
\toprule\noalign{}
t & Remaining product bound \\
\midrule\noalign{}
\endhead
\bottomrule\noalign{}
\endlastfoot
3 & 216 \\
4 & 72 \\
5 & 36 \\
6 & 24 \\
7 & 24 \\
\end{longtable}

Each is below M/9, so one fresh point completes the cover. This excludes
every counterexample with M\textgreater=24697.

\section{8. Finite case with no large prime ---
complete}\label{finite-case-with-no-large-prime-complete}

Every source is 7-smooth. \nolinkurl{smooth/SMOOTH9_PROOF.md} proves this
case independently of common-factor normalization or upward closure.

Its analytic tail starts at M=12555. Below that cutoff, a graph joins
source pairs for which no valid three-bin certificate has been found.
Every possible counterexample is among the resulting seven-source
cliques. The independent enumeration covers 3,413,772 source sets at 94
maxima, with 135,961 distinct demand vectors; the largest surviving
source maximum is 567.

Separate-prime covers, ordinary bins, cumulative-GCD certificates, and
exact CRT profile enumeration close every demand. There are 36,247
explicit slot lists. The final CRT method checks 318 targets over all
relevant prime-power residue classes, not sampled shifts, using 99,864
branch nodes and 6,049 exact multicover calls. Capping valuations at the
source-relevant prime-power level gives lower supply profiles for every
actual shift, so successful covers of all capped profiles prove the
universal result.

The independent replay is \nolinkurl{smooth/verify_smooth9.py}; its
recorded result is \nolinkurl{smooth/verification_result.json}. The CRT
implementation and finite-state bounds were independently audited.

\section{9. Finite case with one large prime ---
complete}\label{finite-case-with-one-large-prime-complete}

The one-large-prime computation covers 59,982 initial configurations.
Universal cofactor envelopes, pair tests, and forced-multiple closure
leave 623 configurations, all with M\textless=735. Exact enumeration
covers 7,006,656 normalized source sets and 658,925 distinct demand
vectors. Their maximal-demand antichain consists of 120,285 demands
closed by independent core, prime-union, or ordinary-bin tests, 252,175
demands with explicit nine-slot certificates, and one hand-proved
demand:

\begin{verbatim}
M=55, demand at (2,3,5,7,11) = (4,9,2,0,4).
\end{verbatim}

For the hand demand, choose two consecutive multiples of 27, a fresh
multiple of 18, and a fresh multiple of 6. The first pair contains an
even number, so these four points supply 2\^{}3*3\^{}9. A multiple of 50
adds 5\^{}2 and a further factor 2; if that multiple was already chosen,
append a fresh even point instead. Thus five points cover
2\textsuperscript{4\emph{3\^{}9}5}2. Unite them with four distinct
multiples of 11. The demands are coprime, so their intersection is
harmless and at most nine points suffice.

\nolinkurl{one/ONE_LARGE_PROOF.md} gives the detailed argument. The
independent source review in \nolinkurl{one/INDEPENDENT_AUDIT.md} checks
the configuration bounds, cofactor envelopes, complete forced-closure
recursion, core-cover conditions, demand encoding, and exact domination
algorithm. The portable \nolinkurl{one/verify_one.py} regenerates all
source configurations, checks every slot certificate using unbounded
integers, and rechecks all 120,285 other maximal demands in Python. Its
independent replay returned PASS. The hand proof removes any dependency
on a CRT computation for the last demand.

\section{10. Finite case with two large primes ---
complete}\label{finite-case-with-two-large-primes-complete}

If no source contains both p and q, their supports are disjoint and one
has size three. If that prime is high, all three pair products of its
cofactors are at most M. If it is low and all its source valuations are
one, the square-gain lemma applies. Otherwise at least two distinct
cofactor-pair products involving a squared-source cofactor are less than
M. Every such product is q-free, while at least three of the four old
sources contain q. Hence at most one old source can block these two
choices, and an unblocked prime-group compression gives nine points.
Thus an actual pq core must occur in a counterexample.

The two-prime finite envelope covers 2,980,414 eligible core classes:
2,784,139 close directly and 196,275 pass to further enumeration.
Pair-graph and closure enumeration on 50,660 configurations gives
112,998 nontrivial-source cliques. Of these, 67,663 close by
core/small-factor criteria; the other 45,335 rows expand to 116,394
actual full source sets, at maxima at most 665.

These sets have 33,939 distinct valuation demands. Componentwise
domination reduces them to 24,511 maximal demands at fixed (p,q,M).
Every maximal demand has a nine-slot certificate. The envelope and
source-enumeration code have independent review. A separate exact
program rebuilt all 116,394 full source sets from the nontrivial rows,
checked normalized source closure, verified equality with all 33,939
stored demand keys, checked that every demand is dominated by a stored
maximal demand, and validated all 24,511 cumulative-GCD certificates.

Primary files are under \nolinkurl{two/}. The independent expansion and
domination check is \nolinkurl{two/verify_demands_independent.py}; it
returned PASS. The complete portable two-prime replay also returned PASS
after rebuilding the finite envelope and enumerator from source.

\section{11. Finite case with three large primes ---
complete}\label{finite-case-with-three-large-primes-complete}

\nolinkurl{three/SUPPORT_THREE.md} eliminates support exactly three for
each large prime in this case. It uses prime-group compression, the
square gain, and a four-atom K4-plus-repeated-edge lemma. Two bounded
residual families contain 162 and 283 source sets. All 445 have exact
nine-slot certificates. The argument explicitly avoids double-counting
coefficients that share primes with an auxiliary composite atom.

Every large prime therefore has source support at least four. Every two
supports intersect, so all three pair cores occur. For
p\textless q\textless r, at most one prime can exceed sqrt(M), and

\begin{verbatim}
qr<=M<9p^2,
\end{verbatim}

with q\^{}2\textless=M. There are 46 possible multiplicity vectors for
the ten cores 1,p,q,r,p\textsuperscript{2,q}2,r\^{}2,pq,pr,qr. The
finite envelope covers 6,294 prime triples, 656,193 eligible maxima, and
850,366 classes.

The envelope uses only proved sufficient criteria: independent
prime-power covers, core banks, spare-bank triangle avoidance, coprime
merging, and exact upper bounds for each rough-prime part of all
possible coefficient choices. It closes 850,364 classes.

The two remaining classes have p=11,q=13,r=23, M=1012 or 1035, and core
pattern

\begin{verbatim}
r,r,pq,pq,pq,pr,qr.
\end{verbatim}

Their full source enumeration has 17,640 raw choices and 1,096
normalized sets, with 347 distinct demands. All 1,096 have universal
nine-slot certificates. An independent program regenerated the actual
sets from integer pools, verified exact equality with the certificate
set, and checked every certificate.

The consolidated proof is \nolinkurl{three/THREE_LARGE_PROOF.md}. Running
\nolinkurl{three/verify_all_three.py} regenerates both support-three
finite families, verifies their 445 certificates, recompiles and reruns
all 850,366 envelope classes in an isolated directory, checks the exact
residual classes, and independently regenerates and checks all 1,096
final certificates. The complete replay returned PASS. It needs only
Python's standard library and a C++17 compiler.

\section{12. Conclusion and limits}\label{conclusion-and-limits}

The case split is exhaustive. A normalized counterexample can have only
zero, one, two, or three large primes; four have been excluded
analytically. The global analytic argument removes every
M\textgreater=24697. Each finite branch below that threshold has now
been closed, regenerated, and independently checked. Therefore no
counterexample exists, and g(7)\textless=9. The explicit example in
Section 3 proves g(7)\textgreater=9. Together these establish

\begin{verbatim}
g(7)=9.
\end{verbatim}

This is a computer-assisted mathematical proof, with the stated prior
small-n inputs. Its finite source ranges and finite CRT classes become
universal only through the explicit reductions above. The cumulative-GCD
certificates prove covers for every interval shift; the CRT certificates
exhaust all relevant capped prime-power profiles. No conclusion is based
on testing a finite list of numerical shifts.

This proof establishes no new exact value for g(8) or g(9), and it does
not settle either global coefficient-two conjecture in Erdős Problem
708.

\end{document}
